\section{Introduction}

\IEEEPARstart{T}{he} idea of a self-replicating space factory is old and simple to state:
land one factory on the Moon, an asteroid, or Mars, and let it build copies of itself from
local material, so that an entire industry grows from a single rocket's worth of cargo. It
begins with von Neumann, who proved that a machine can build a copy of itself if it carries
both a description of its construction and a way to copy that description
\cite{von-neumann-burks-1966}. The canonical engineering study is NASA's seed-factory
design, which framed the concept in terms of ``vitamins'': the factory makes almost all of
itself locally and imports only the parts it cannot make \cite{nasa-cp-2255-1980}. The
modern synthesis restated the same idea as \emph{closure}, the degree to which a machine
can reproduce its own matter, energy, and information, and identified the dynamic that
matters most: below full closure the factory waits on resupply, and only near full closure
does its growth become self-sustaining \cite{freitas-merkle-2004}.

The concept is neither purely theoretical nor purely futuristic. Concrete lunar-factory
architectures decompose the machine into mining, refining, parts fabrication, and assembly
subsystems \cite{chirikjian-2004-niac,moses-chirikjian-2020}; physical prototypes have
demonstrated mechanical self-reproduction from supplied modules \cite{zykov-lipson-2005};
and the open-source RepRap 3D printer prints a large fraction of its own parts while
importing its motors, rods, and electronics \cite{reprap-jones-2011}, a terrestrial echo of
the wall this paper is about. Recent quantitative studies of space manufacturing continue
the pattern: a tonne-scale landed seed bootstrapping to thousands of tonnes of industrial
assets over about two decades \cite{metzger-2013}, a self-replicating interstellar probe
whose replicated mass fraction is about 70 percent with electronics imported
\cite{borgue-hein-2020}, and a colony-scale factory that likewise keeps chips and circuitry
Earth-sourced \cite{guided-self-replicating-factory-2021}. Across every one of these, the
same part is left off the list of what the machine can make.

This paper isolates one question inside that framing. If closure is the whole game, what
stops it from reaching one? The answer, consistently, is chips. A lone factory can smelt
metal, cast structure, and assemble mechanisms, but it cannot make its own electronics. We
argue that this is not an incidental gap to be closed later but a wall with two independent
foundations, supply-chain depth and embodied energy, and that the wall is exactly what
bounds the launch-mass leverage that motivates self-replication in the first place.
